\input{TAC-801-SHARED-PREAMBLE.tex}
\begin{document}
\nolinenumbers
\begin{singlespace}
Franciscus Dylan Rosario, \\
7624 Melody \\
Rohnert Park, CA 94928 \\
E: soltrinox@gmail.com \\
Plaintiff, In Pro Per \\
\end{singlespace}
\vspace*{9mm}
\begin{tightcenter}
\textbf{SUPERIOR COURT OF THE STATE OF CALIFORNIA} \\
\textbf{COUNTY OF SAN FRANCISCO}
\end{tightcenter}
\vspace*{2.25mm}
\nolinenumbers
\begin{minipage}[t]{3in}
    \raggedright
    \textbf{Franciscus Dylan Rosario,} \\ \textbf{PLAINTIFF}, \\
    \hspace{1cm} v. \\
    \small
    CSAA Insurance Exchange; \\
    Carbone, Smith \& Koyama, LLP; \\
    Michael R. Chambers; \\
    Phillips, Spallas \& Angstadt LLP; \\
    Priya D. Navaratnasingham; \\
    Alberto Reyna; \\
    and DOES 1 through 20, inclusive; \\
    \normalsize
    \textbf{DEFENDANTS}
    \makebox[3in]{\hrulefill}
\end{minipage}%
\hspace{3mm}
\begin{minipage}[t]{0.5pt}
    \vspace{0pt}
    \rule{0.5pt}{3.1in}
\end{minipage}
\hspace{3mm}
\begin{minipage}[t]{3in}
    \raggedright
    \noindent \textbf{Case No.: CGC-25-631801} \\
    \textbf{(Related: CGC-21-594102)} \\
    ~\\
    \textbf{CURE MATRIX — DEFENSE OBJECTION MAP}
\end{minipage}
\vspace*{6mm}
\linenumbers
\section{Second-Layer Cure Matrix If SAC Defects Survive}\label{second-layer-cure-matrix-if-sac-defects-survive}

\textbf{Case No.:} CGC-25-631801\\
\textbf{Subject:} Third Amended Complaint for Extrinsic Fraud\\
\textbf{Status:} CONTINGENT --- filed only if Round D defects surface after SAC ruling\\
\textbf{Purpose:} Maps each defense objection from the April 22, 2026 CCP \textsection{} 436 Motion to Strike to the specific TAC paragraph(s) that cure it.

\textbf{First Cure:} SAC-leave motion (hearing May 6, 2026)\\
\textbf{Second Cure:} This TAC packet (hearing TBD on OST if needed)

\textbf{Underlying Action:} Rosario v. Abdelhalim, Case No.~CGC-21-594102

\begin{center}\rule{0.5\linewidth}{0.5pt}\end{center}

\section{Master Cure Matrix --- Nine Defense Objections}\label{master-cure-matrix-nine-defense-objections}

{\def\LTcaptype{table} % do not increment counter
\begin{longtable}[]{@{}
  >{\raggedright\arraybackslash}p{(\linewidth - 14\tabcolsep) * \real{0.0214}}
  >{\raggedright\arraybackslash}p{(\linewidth - 14\tabcolsep) * \real{0.1357}}
  >{\raggedright\arraybackslash}p{(\linewidth - 14\tabcolsep) * \real{0.1357}}
  >{\raggedright\arraybackslash}p{(\linewidth - 14\tabcolsep) * \real{0.0714}}
  >{\raggedright\arraybackslash}p{(\linewidth - 14\tabcolsep) * \real{0.1143}}
  >{\raggedright\arraybackslash}p{(\linewidth - 14\tabcolsep) * \real{0.1071}}
  >{\raggedright\arraybackslash}p{(\linewidth - 14\tabcolsep) * \real{0.1929}}
  >{\raggedright\arraybackslash}p{(\linewidth - 14\tabcolsep) * \real{0.2214}}@{}}
\toprule\noalign{}
\begin{minipage}[b]{\linewidth}\raggedright
\#
\end{minipage} & \begin{minipage}[b]{\linewidth}\raggedright
Defense Objection
\end{minipage} & \begin{minipage}[b]{\linewidth}\raggedright
Defense MPA Basis
\end{minipage} & \begin{minipage}[b]{\linewidth}\raggedright
TAC Cure
\end{minipage} & \begin{minipage}[b]{\linewidth}\raggedright
TAC Section/¶¶
\end{minipage} & \begin{minipage}[b]{\linewidth}\raggedright
Key Authority
\end{minipage} & \begin{minipage}[b]{\linewidth}\raggedright
Cal. Supreme Court Anchor
\end{minipage} & \begin{minipage}[b]{\linewidth}\raggedright
Bench Brief/SAC Back-Reference
\end{minipage} \\
\midrule\noalign{}
\endhead
\bottomrule\noalign{}
\endlastfoot
1 & \emph{Kulchar} extrinsic-fraud not viable & Independent equity action improper & Four-channel \emph{Kulchar} framework; judicial-reliance quote; seven consequential rulings; Supreme Court spine \textsection{}I.4.5 & \textsection{}I, ¶¶ 1--7, 7.5; \textsection{}I.4.5, ¶¶ 3.5--3.7; \textsection{}I.5, ¶¶ 4--6 & \emph{Kougasian v. TMSL} (2004) 118 Cal.App.4th 1028 & \emph{Kulchar v. Kulchar} (1969) 1 Cal.3d 467, 471; \emph{In re Marriage of Modnick} (1983) 33 Cal.3d 897, 905; \emph{Estate of Sanders} (1985) 40 Cal.3d 607, 614 & Bench Brief \textsection{}\textsection{}I.A-I.C \\
2 & Appeal A173827 is the remedy & CCP \textsection{} 904.1 routing & \emph{Kougasian/Stevenot} independent action doctrine; \emph{Throckmorton} pre-rebuttal & \textsection{}IV.5, ¶¶ 28--31 & \emph{In re Marriage of Stevenot} (1984) 154 Cal.App.3d 1051 & \emph{Estate of Sanders} (1985) 40 Cal.3d 607, 614 (independent action in equity recognized by Supreme Court) & Bench Brief Part VII \\
3 & Civil Code \textsection{} 47(b) privilege bars claims & Litigation privilege & Nine-Point Evidentiary Chain; non-communicative conduct; \emph{Flatley} illegality anchor; \emph{Action Apartment} extrinsic-fraud exception & \textsection{}II, ¶¶ 8--13 & \emph{Flatley v. Mauro} (2006) 39 Cal.4th 299; \emph{Action Apartment} (2007) 41 Cal.4th 1232 & \emph{Estate of Sanders} (1985) 40 Cal.3d 607, 614 (equity exception); \emph{In re Marriage of Modnick} (1983) 33 Cal.3d 897, 907 (concealment as fraud) & SAC-06 \textsection{}IV.D.0 \\
4 & Standing --- \emph{Shaolian/Norton} & Attorney-client privity & \emph{Estate of Sanders} distinction; new ¶ 15.5 officer-of-court fiduciary-duty theory; \emph{Doctors' Company} aiding-and-abetting framework; Civ. Code \textsection{} 2338 agency/ratification & \textsection{}III, ¶¶ 14--21, 15.5 & \emph{Estate of Sanders} (1985) 40 Cal.3d 607; \emph{Doctors' Company} (1989) 49 Cal.3d 39 & \emph{Estate of Sanders} (1985) 40 Cal.3d 607, 615 (fiduciary-concealment holding) & Bench Brief \textsection{}II.F; SAC-06 ¶¶ 7--14 \\
5 & Collateral estoppel / res judicata & Underlying judgment preclusive & Five-issues-not-adjudicated block; primary-rights comparison table; \emph{Mycogen/Boeken} same-primary-right test; new ¶ 26.5 four-Supreme-Court recital & \textsection{}IV, ¶¶ 22--27, 26.5 & \emph{Lucido v. Superior Court} (1990) 51 Cal.3d 335; \emph{Mycogen} (2002) 28 Cal.4th 888 & \emph{Mycogen}, 28 Cal.4th at 904; \emph{Boeken v. Philip Morris USA} (2010) 48 Cal.4th 788, 798; \emph{Crowley v. Katleman} (1994) 8 Cal.4th 666, 681--682; \emph{Lucido}, 51 Cal.3d at 341 & SAC-06 \textsection{}IV.B \\
6 & CCP \textsection{} 1032 prevailing-party overlap & Costs already adjudicated & Express disclaimer (\textsection{}V, ¶ 32(b)) --- no recovery of prevailing-party costs except via vacatur & \textsection{}V, ¶ 32(b) & --- & Same primary-rights anchors as Objection 5 (no standalone Supreme Court anchor required) & --- \\
7 & Damages duplicative of PI case & Personal injury damages re-litigated & Reservation paragraph; express disclaimers; fraud-delta limitation; new \textsection{}IV.5.5 ``Opportunity to Present, Not Admission'' & \textsection{}V, ¶¶ 32--34; \textsection{}IV.5.5, ¶ 34.7 & --- & \emph{Westphal v. Westphal} (1942) 20 Cal.2d 393, 397; \emph{Modnick}, 33 Cal.3d at 907 (deprivation, not exhibit count) & --- \\
8 & IIED / \emph{Ribas} claim improper & Standalone emotional distress & No standalone IIED claim asserted; emotional distress pleaded as equity factor only & \textsection{}VI, Tier 4, ¶ 38 & --- & \emph{Estate of Sanders} (1985) 40 Cal.3d 607, 614 (equity-factor frame; no standalone IIED) & --- \\
9 & Ultimate admission cures fraud (``it got in on Day 7'') & Late admission moots extrinsic-fraud injury & New \textsection{}IV.5.5 and ¶¶ 34.5--34.6 measuring deprivation at opportunity-to-present, not at admission; Memo \textsection{} V.G & TAC \textsection{}IV.5.5, ¶ 34.7; ¶¶ 34.5--34.6; Memo \textsection{} V.G & --- & \emph{Westphal}, 20 Cal.2d at 397; \emph{Kulchar}, 1 Cal.3d at 471; \emph{Modnick}, 33 Cal.3d at 907; \emph{Sanders}, 40 Cal.3d at 616 (quoting \emph{Larrabee v. Tracy} (1943) 21 Cal.2d 645, 650--651) & Memo \textsection{} I.F; TAC Eighth Cause of Action \\
\end{longtable}
}

\begin{center}\rule{0.5\linewidth}{0.5pt}\end{center}

\section{Detailed Cure Analysis}\label{detailed-cure-analysis}

\subsection{Objection 1: Kulchar Extrinsic-Fraud Framework Not Pleaded}\label{objection-1-kulchar-extrinsic-fraud-framework-not-pleaded}

\textbf{Defense Argument:} The complaint fails to plead the elements of an independent action in equity under \emph{Kulchar v. Kulchar} (1969) 1 Cal.3d 467.

\textbf{TAC Cure:}

\begin{itemize}
\tightlist
\item
  \textbf{¶ 1:} Identifies action as independent action in equity under \emph{Kougasian/Olivera/Stevenot}
\item
  \textbf{¶ 2:} Distinguishes from appeal, motion for new trial, or collateral attack
\item
  \textbf{¶ 3:} Four-channel \emph{Kulchar} framework pleaded element-by-element:

  \begin{itemize}
  \tightlist
  \item
    \begin{enumerate}
    \def\labelenumi{(\alph{enumi})}
    \tightlist
    \item
      Concealment of facts (production history)
    \end{enumerate}
  \item
    \begin{enumerate}
    \def\labelenumi{(\alph{enumi})}
    \setcounter{enumi}{1}
    \tightlist
    \item
      False promises/representations of non-receipt
    \end{enumerate}
  \item
    \begin{enumerate}
    \def\labelenumi{(\alph{enumi})}
    \setcounter{enumi}{2}
    \tightlist
    \item
      Keeping Plaintiff in ignorance (893-day silence)
    \end{enumerate}
  \item
    \begin{enumerate}
    \def\labelenumi{(\alph{enumi})}
    \setcounter{enumi}{3}
    \tightlist
    \item
      Misleading the tribunal (categorical denials)
    \end{enumerate}
  \end{itemize}
\item
  \textbf{¶¶ 4--6:} Judicial-reliance quote and seven consequential rulings
\end{itemize}

\textbf{Authority:} \emph{Kougasian v. TMSL, Inc.} (2004) 118 Cal.App.4th 1028, 1032--35

\textbf{Cal. Supreme Court Anchor:} \emph{Kulchar v. Kulchar} (1969) 1 Cal.3d 467, 471 (``fair adversary hearing'' primary right); \emph{In re Marriage of Modnick} (1983) 33 Cal.3d 897, 905 (broad-concept extrinsic fraud); \emph{Estate of Sanders} (1985) 40 Cal.3d 607, 614 (fiduciary-breach conversion to extrinsic). TAC installs these anchors at new \textsection{}I.4.5 (¶¶ 3.5--3.7) and new ¶ 7.5.

\textbf{Back-Reference:} Bench Brief \textsection{}\textsection{} I.A-I.C (judicial reliance sequence)

\begin{center}\rule{0.5\linewidth}{0.5pt}\end{center}

\subsection{Objection 2: Appeal A173827 Is the Remedy}\label{objection-2-appeal-a173827-is-the-remedy}

\textbf{Defense Argument:} Plaintiff's remedy is the pending appeal; no independent action is cognizable.

\textbf{TAC Cure:}

\begin{itemize}
\tightlist
\item
  \textbf{¶¶ 28--29:} Distinguishes appeal (instructional error) from this action (institutional fraud)
\item
  \textbf{¶ 30:} Cites \emph{Kougasian} and \emph{Stevenot} for parallel-remedy principle
\item
  \textbf{¶ 31:} Pre-rebuts \emph{Throckmorton} as non-binding federal authority
\end{itemize}

\textbf{Authority:} \emph{In re Marriage of Stevenot} (1984) 154 Cal.App.3d 1051; \emph{Kougasian}, 118 Cal.App.4th at 1034--36

\begin{center}\rule{0.5\linewidth}{0.5pt}\end{center}

\subsection{Objection 3: Litigation Privilege Bars Claims}\label{objection-3-litigation-privilege-bars-claims}

\textbf{Defense Argument:} All claims are barred by the absolute litigation privilege under Civil Code \textsection{} 47(b).

\textbf{TAC Cure:}

\begin{itemize}
\tightlist
\item
  \textbf{¶ 8:} Identifies conduct as non-communicative (silence, file-transfer, failure to correct)
\item
  \textbf{¶ 9:} Nine-Point Evidentiary Chain documenting 893 days of defense possession
\item
  \textbf{¶¶ 10--11:} \emph{Flatley} illegality exception (Bus. \& Prof.~Code \textsection{} 6128(a))
\item
  \textbf{¶ 12:} \emph{Action Apartment} extrinsic-fraud carveout
\item
  \textbf{¶ 13:} \emph{Throckmorton} pre-rebuttal
\end{itemize}

\textbf{Authority:} - \emph{Rusheen v. Cohen} (2006) 37 Cal.4th 1048 (non-communicative conduct) - \emph{Flatley v. Mauro} (2006) 39 Cal.4th 299, 320--325 (illegality exception) - \emph{Action Apartment Ass'n v. City of Santa Monica} (2007) 41 Cal.4th 1232, 1248--49 (extrinsic-fraud carveout)

\textbf{Cal. Supreme Court Anchor:} \emph{Estate of Sanders} (1985) 40 Cal.3d 607, 614 (equity exception for officer-of-court fraud); \emph{In re Marriage of Modnick} (1983) 33 Cal.3d 897, 907 (denial of existence as extrinsic concealment).

\textbf{Back-Reference:} SAC-06 \textsection{}IV.D.0 (Nine-Point Chain); Bench Brief \textsection{}IV.E

\begin{center}\rule{0.5\linewidth}{0.5pt}\end{center}

\subsection{Objection 4: Lack of Standing (Shaolian/Norton)}\label{objection-4-lack-of-standing-shaoliannorton}

\textbf{Defense Argument:} Plaintiff lacks standing to sue defense attorneys due to attorney-client privity requirements.

\textbf{TAC Cure:}

\begin{itemize}
\tightlist
\item
  \textbf{¶ 15:} \emph{Shaolian} distinction (policy benefits vs.~fraud-on-court)
\item
  \textbf{¶ 15(c):} \emph{Estate of Sanders} (fraud-on-court parties)
\item
  \textbf{¶¶ 16--17:} Civ. Code \textsection{} 2338 agency/ratification
\item
  \textbf{¶ 18:} \emph{Doctors' Company} aiding-and-abetting framework
\item
  \textbf{¶ 19:} \textsection{} 6128(a) independent tort under \emph{Rickley/Shaffer}
\item
  \textbf{¶¶ 20--21:} Civ. Code \textsection{} 1573 constructive fraud
\end{itemize}

\textbf{Authority:} - \emph{Estate of Sanders} (1985) 40 Cal.3d 607, 615--616 - \emph{Doctors' Company v. Superior Court} (1989) 49 Cal.3d 39 - \emph{Rickley v. Goodfriend} (2013) 212 Cal.App.4th 1136

\textbf{Cal. Supreme Court Anchor:} \emph{Estate of Sanders} (1985) 40 Cal.3d 607, 615 (officer-of-court fiduciary duty runs to adjudicative integrity; opposing party has equity standing regardless of privity). TAC installs this anchor at new ¶ 15.5 and at new Eighth Cause of Action.

\textbf{Back-Reference:} Bench Brief \textsection{}II.F; SAC-06 ¶¶ 7--14 (party block)

\begin{center}\rule{0.5\linewidth}{0.5pt}\end{center}

\subsection{Objection 5: Collateral Estoppel / Res Judicata}\label{objection-5-collateral-estoppel-res-judicata}

\textbf{Defense Argument:} Action is a collateral attack on a valid judgment barred by res judicata.

\textbf{TAC Cure:}

\begin{itemize}
\tightlist
\item
  \textbf{¶ 22:} Distinct primary right: freedom from fraud on the court
\item
  \textbf{¶ 24:} Five-issues-not-adjudicated block
\item
  \textbf{¶ 25:} Primary-rights comparison table
\item
  \textbf{¶¶ 26--27:} \emph{Mycogen/Crowley/Lucido} recital
\end{itemize}

\textbf{Authority:} - \emph{Lucido v. Superior Court} (1990) 51 Cal.3d 335 - \emph{Mycogen Corp.~v. Monsanto Co.} (2002) 28 Cal.4th 888, 904 - \emph{Boeken v. Philip Morris USA, Inc.} (2010) 48 Cal.4th 788, 798

\textbf{Cal. Supreme Court Anchor:} \emph{Mycogen}, 28 Cal.4th at 904 (primary right = injury suffered); \emph{Boeken}, 48 Cal.4th at 798 (primary right independent of forum); \emph{Crowley v. Katleman} (1994) 8 Cal.4th 666, 681--682 (separate injuries {\textrightarrow{}} separate primary rights); \emph{Lucido}, 51 Cal.3d at 341 (actually-litigated + necessarily-decided). TAC installs these at new ¶ 26.5.

\textbf{Back-Reference:} SAC-06 \textsection{}IV.B (Issues Not Adjudicated)

\begin{center}\rule{0.5\linewidth}{0.5pt}\end{center}

\subsection{Objection 6: CCP \textsection{} 1032 Prevailing-Party Overlap}\label{objection-6-ccp-1032-prevailing-party-overlap}

\textbf{Defense Argument:} Plaintiff seeks to recover costs already adjudicated under \textsection{} 1032.

\textbf{TAC Cure:}

\begin{itemize}
\tightlist
\item
  \textbf{¶ 32(b):} Express disclaimer --- Plaintiff does not seek to recover prevailing-party costs from underlying trial, except to the extent such costs would be vacated as a consequence of vacating the underlying judgment.
\end{itemize}

\begin{center}\rule{0.5\linewidth}{0.5pt}\end{center}

\subsection{Objection 7: Damages Duplicative of PI Case}\label{objection-7-damages-duplicative-of-pi-case}

\textbf{Defense Argument:} Plaintiff seeks to recover PI damages already litigated in CGC-21-594102.

\textbf{TAC Cure:}

\begin{itemize}
\tightlist
\item
  \textbf{¶ 32:} Express disclaimers:

  \begin{itemize}
  \tightlist
  \item
    \begin{enumerate}
    \def\labelenumi{(\alph{enumi})}
    \tightlist
    \item
      No PI damages from underlying action
    \end{enumerate}
  \item
    \begin{enumerate}
    \def\labelenumi{(\alph{enumi})}
    \setcounter{enumi}{1}
    \tightlist
    \item
      No prevailing-party costs (except via vacatur)
    \end{enumerate}
  \item
    \begin{enumerate}
    \def\labelenumi{(\alph{enumi})}
    \setcounter{enumi}{2}
    \tightlist
    \item
      No duplicative 802 damages
    \end{enumerate}
  \item
    \begin{enumerate}
    \def\labelenumi{(\alph{enumi})}
    \setcounter{enumi}{3}
    \tightlist
    \item
      No pro-per attorney's fees
    \end{enumerate}
  \end{itemize}
\item
  \textbf{¶¶ 33--34:} Partial-presentation rebuttal
\item
  \textbf{\textsection{}IV.5.5, ¶ 34.7:} Deprivation measured by opportunity to present, not admission
\item
  \textbf{\textsection{}VI, Tier 2, ¶ 36:} Fraud-delta limitation on compensatory relief
\end{itemize}

\textbf{Cal. Supreme Court Anchor:} \emph{Westphal v. Westphal} (1942) 20 Cal.2d 393, 397; \emph{In re Marriage of Modnick} (1983) 33 Cal.3d 897, 907 (measure is deprivation of opportunity, not exhibit count).

\begin{center}\rule{0.5\linewidth}{0.5pt}\end{center}

\subsection{Objection 8: IIED / Ribas Claim Improper}\label{objection-8-iied-ribas-claim-improper}

\textbf{Defense Argument:} Plaintiff improperly asserts a standalone IIED claim barred by \emph{Ribas} or \textsection{} 47(b).

\textbf{TAC Cure:}

\begin{itemize}
\tightlist
\item
  \textbf{¶ 38:} No standalone IIED cause of action asserted
\item
  Emotional distress pleaded only as Tier 4 equity factor supporting prejudice, not as independent tort
\end{itemize}

\begin{center}\rule{0.5\linewidth}{0.5pt}\end{center}

\subsection{Objection 9: Ultimate Admission Cures Fraud (``It Got In On Day 7'')}\label{objection-9-ultimate-admission-cures-fraud-it-got-in-on-day-7}

\textbf{Defense Argument (Anticipated):} The Day-7 admission of the 911 audio, after Defendant Abdelhalim authenticated it as his own voice, retroactively cures any prior concealment and moots the extrinsic-fraud injury.

\textbf{TAC Cure:}

\begin{itemize}
\tightlist
\item
  \textbf{\textsection{}IV.5.5, ¶ 34.7:} Installs the Supreme Court deprivation rule --- \emph{Westphal}, \emph{Kulchar}, \emph{Modnick}, \emph{Sanders} --- as a four-test foreclosure of the Day-7 defense.
\item
  \textbf{¶ 34.5:} Applies \emph{Modnick} 907 ``fairly before the court for adjudication'' standard to the seven rulings issued within the corrupted record.
\item
  \textbf{¶ 34.6:} Applies \emph{Sanders} 616 / \emph{Larrabee} 650--651 ``effectively prevented'' standard to the Court's decision-making matrix.
\item
  \textbf{¶ 33 (edited):} Embeds \emph{Westphal} 397 ``opportunity to present'' and \emph{Kulchar} 471 ``fair adversary hearing'' into the partial-presentation rebuttal.
\item
  \textbf{Eighth Cause of Action (¶¶ 81--87):} Officer-of-the-court fiduciary-duty theory anchored in \emph{Sanders} 615.
\item
  \textbf{Memo \textsection{} V.G:} Parallel argument in the memo.
\end{itemize}

\textbf{Authority (Cal. Supreme Court Only):}

\begin{itemize}
\tightlist
\item
  \emph{Westphal v. Westphal} (1942) 20 Cal.2d 393, 397 (extrinsic = ``prevented from exhibiting fully his case'')
\item
  \emph{Kulchar v. Kulchar} (1969) 1 Cal.3d 467, 471 (primary right = ``fair adversary hearing'')
\item
  \emph{In re Marriage of Modnick} (1983) 33 Cal.3d 897, 907 (``never fairly before the court for adjudication'')
\item
  \emph{Estate of Sanders} (1985) 40 Cal.3d 607, 616 (conduct that ``effectively prevented \ldots{} a fair adversary hearing''), quoting \emph{Larrabee v. Tracy} (1943) 21 Cal.2d 645, 650--651
\end{itemize}

\textbf{Cal. Supreme Court Anchor:} All four holdings above. The Supreme Court has located the operative injury at the point of deprivation of opportunity, not at the point of eventual admission. The Day-7 leak-through of a single exhibit does not retroactively cure the seven structural rulings issued within the corrupted \emph{Modnick} record.

\textbf{Back-Reference:} Memo \textsection{} I.F; Memo \textsection{} V.G; TAC Eighth Cause of Action.

\begin{center}\rule{0.5\linewidth}{0.5pt}\end{center}

\section{Cross-Reference: TAC Sections to Defense Objections}\label{cross-reference-tac-sections-to-defense-objections}

{\def\LTcaptype{table} % do not increment counter
\begin{longtable}[]{@{}
  >{\raggedright\arraybackslash}p{(\linewidth - 4\tabcolsep) * \real{0.3421}}
  >{\raggedright\arraybackslash}p{(\linewidth - 4\tabcolsep) * \real{0.1842}}
  >{\raggedright\arraybackslash}p{(\linewidth - 4\tabcolsep) * \real{0.4737}}@{}}
\toprule\noalign{}
\begin{minipage}[b]{\linewidth}\raggedright
TAC Section
\end{minipage} & \begin{minipage}[b]{\linewidth}\raggedright
Title
\end{minipage} & \begin{minipage}[b]{\linewidth}\raggedright
Cures Objections
\end{minipage} \\
\midrule\noalign{}
\endhead
\bottomrule\noalign{}
\endlastfoot
I & Character of the Action & 1 \\
I.4.5 & Supreme Court--Recognized Right to a Fair Adversary Hearing (¶¶ 3.5--3.7) & 1, 9 \\
I.5 & Court's Articulated Reliance & 1 \\
II & Non-Communicative Conduct Block & 3 \\
III & Standing Block (incl.~¶ 15.5 officer-of-court fiduciary duty) & 4 \\
IV & Primary Rights Statement (incl.~¶ 26.5 four-Supreme-Court recital) & 5 \\
IV.5 & Appeal Non-Preclusion & 2 \\
IV.5.5 & Measure of Deprivation: Opportunity to Present, Not Admission (¶ 34.7) & 7, 9 \\
V & Reservation Paragraph (incl.~¶¶ 34.5--34.6 \emph{Modnick}/\emph{Sanders} standards) & 6, 7, 9 \\
VI & Tiered Prayer for Relief & 8 \\
VII & Party-Specific Allegations & --- \\
VIII & Day-7 Admission Foundation & 1, 3, 9 \\
IX & Laws-Broken Catalog & 3 \\
Eighth CoA & Extrinsic Fraud by Officer-of-Court Fiduciary Breach (¶¶ 81--87) & 4, 9 \\
\end{longtable}
}

\begin{center}\rule{0.5\linewidth}{0.5pt}\end{center}

\section{Uncured Gaps (Acknowledged for Strategic Purposes)}\label{uncured-gaps-acknowledged-for-strategic-purposes}

The following items are \textbf{not} cured by the TAC and are acknowledged for strategic transparency:

{\def\LTcaptype{table} % do not increment counter
\begin{longtable}[]{@{}
  >{\raggedright\arraybackslash}p{(\linewidth - 2\tabcolsep) * \real{0.2174}}
  >{\raggedright\arraybackslash}p{(\linewidth - 2\tabcolsep) * \real{0.7826}}@{}}
\toprule\noalign{}
\begin{minipage}[b]{\linewidth}\raggedright
Gap
\end{minipage} & \begin{minipage}[b]{\linewidth}\raggedright
Reason Not Cured
\end{minipage} \\
\midrule\noalign{}
\endhead
\bottomrule\noalign{}
\endlastfoot
Anti-SLAPP cross-briefing & Handled separately; TAC notes CCP \textsection{} 425.17 exemptions only \\
CSAA ``did not actually litigate'' imputation defense & Addressed via Civ. Code \textsection{} 2338 ratification; CSAA funded, directed, and approved litigation posture \\
\end{longtable}
}

\begin{center}\rule{0.5\linewidth}{0.5pt}\end{center}

\section{Validation Checklist}\label{validation-checklist}

\begin{itemize}
\tightlist
\item[$\boxtimes$]
  All eight defense objections mapped to TAC cure
\item[$\boxtimes$]
  Each cure supported by legal authority and/or factual specificity
\item[$\boxtimes$]
  Day-7 admission attributed to Navaratnasingham (RT 7 p.~213:13--15)
\item[$\boxtimes$]
  Judicial-reliance quote included (Hr'g Tr. Feb.~28, 2025, pp.~291:26--292:7)
\item[$\boxtimes$]
  Nine-Point Evidentiary Chain included
\item[$\boxtimes$]
  Party block matches SAC-06 ¶¶ 7--14 (with Bar Nos.)
\item[$\boxtimes$]
  Laws-broken catalog included
\end{itemize}

\begin{center}\rule{0.5\linewidth}{0.5pt}\end{center}

\emph{Last updated: April 22, 2026}

\nolinenumbers
\input{TAC-801-SIGNATURE-BLOCK.tex}

\end{document}
